\chapter{Allgemeiner Nulltransport}
\hypertarget{zero-transport.proofs}{}

Diese Ergänzung beweist den Nulltransport ohne Endlichkeitsvoraussetzung.
Die \href{_B28-zero-transport-reading.pdf\#zero-transport.reading}{Lesefassung}
erklärt dieselbe Ableitung. Verwendet werden die Mengen- und
Funktionsgrundlagen, die Potenzhalbgruppen aus Band~28 und die
Nullnotation aus Band~33. Alle zusätzlichen Sätze werden hier unter
der eigenen Bandkennung~28Z bewiesen.

Im ganzen Kapitel gelten die Abkürzungen
\[
 \begin{gathered}
 K_S=\PowNE S,\quad K_T=\PowNE T,\quad z_S=\{0_S\},\\
 G_S=\SemigroupAlg{S}{\star},\quad G_T=\SemigroupAlg{T}{\diamond},\quad
 H_S=\SemigroupWithZeroAlg{S}{\star}{0_S},\\
 I=\SemigroupIso{\Phi}{K_S,\star_{\mathcal P}}{K_T,\diamond_{\mathcal P}},\\
 Z=\Phi(z_S),\qquad A_t=\Phi^{-1}(\{t\}),\qquad
 W=\bigcup_{t\in Z}A_t.
 \end{gathered}
\]
Bei Elementen und bei Mengenprodukten lassen wir das durch den Träger
festgelegte Operationszeichen weg. \(A_t\) wird nur für \(t\in T\)
verwendet, sobald \(I\) vorausgesetzt ist. Die indizierte Vereinigung
ist die Schreibweise aus
\FormulaRefAuto{SemigroupIndexedUnionNotationDef}[def].

\clearpage
\section{Vereinigungen und Mengenprodukte}

\Needspace{12\baselineskip}
\FormulaThmDeltaK[Ein Glied liegt in der Vereinigung]{%
 i\in J\vdash B_i\subseteq\bigcup_{j\in J}B_j
}{ZeroTransportUnionContainsMember}{%
 \DeltaRow{Mengenfamilie}{(B_j)_{j\in J}}
 \DeltaRow{Mengen}{J\dsep i\dsep j\dsep x}}
\begin{tabproofwide}
 \proofstepwidestar[1]{i\in J}{\rA}
 \proofstepwidestar[2]{x\in B_i}{\rA}
 \proofstepwidestar[1,2]{\exists j\in J\;(x\in B_j)}{\rEI{\rAI{1,2}}}
 \proofstepwidestar[1,2]{x\in\bigcup_{j\in J}B_j}{\FormulaRefAuto{SemigroupIndexedUnionNotationDef}[def]{3}}
 \proofstepwidestar[1]{B_i\subseteq\bigcup_{j\in J}B_j}{\FormulaRefAuto{A\subseteq B\coloneq\forall x\,(x\in A\rightarrow x\in B)}[def]{\rUI{\rRI{2,4}}}}
\end{tabproofwide}

\Needspace{12\baselineskip}
\FormulaThmDeltaK[Nichtleere Vereinigung nichtleerer Teilmengen]{%
 J\neq\varnothing\dsep\forall j\in J\;(B_j\in\PowNE S)
 \vdash\bigcup_{j\in J}B_j\in\PowNE S
}{ZeroTransportUnionCarrier}{%
 \DeltaRow{Mengenfamilie}{(B_j)_{j\in J}}
 \DeltaRow{Mengen}{S\dsep J\dsep j\dsep x}}
Für diesen Beweis sei \(U=\bigcup_{j\in J}B_j\).
\begin{tabproofsplitwide}
\proofpartwide{Teilmengenzugehörigkeit}
 \proofstepwidestar[1]{J\neq\varnothing}{\rA}
 \proofstepwidestar[2]{\forall j\in J\;(B_j\in\PowNE S)}{\rA}
 \proofstepwidestar[3]{x\in U}{\rA}
 \proofstepwidestar[3]{\exists j\in J\;(x\in B_j)}{\FormulaRefAuto{SemigroupIndexedUnionNotationDef}[def]{3}}
 \proofstepwidestar[5]{j\in J\land x\in B_j}{\rA}
 \proofstepwidestar[2,5]{B_j\in\PowNE S}{\rRE{\rUE{2},\rAEa{5}}}
 \proofstepwidestar[2,5]{x\in S}{\FormulaRefAuto{PowerSemigroupCarrierElement}[thm]{6,\rAEb{5}}}
 \proofstepwidestar[2,3]{x\in S}{\rEE{4,5,7}}
 \proofstepwidestar[2]{U\subseteq S}{\FormulaRefAuto{A\subseteq B\coloneq\forall x\,(x\in A\rightarrow x\in B)}[def]{\rUI{\rRI{3,8}}}}
\closeproofpartwide
\Needspace{8\baselineskip}
\proofpartwide{Nichtleerheit}
 \setcounter{proofstepnr}{9}
 \proofstepwidestar[1]{\exists j\;(j\in J)}{\FormulaRefAuto{A\neq\varnothing\eqvdash\exists x\,(x\in A)}[thm]{1}}
 \proofstepwidestar[11]{j\in J}{\rA}
 \proofstepwidestar[2,11]{B_j\in\PowNE S}{\rRE{\rUE{2},11}}
 \proofstepwidestar[2,11]{B_j\subseteq S\land B_j\neq\varnothing}{\FormulaRefAuto{NonemptyPowersetMembership}[thm]{12}}
 \proofstepwidestar[2,11]{\exists x\;(x\in B_j)}{\FormulaRefAuto{A\neq\varnothing\eqvdash\exists x\,(x\in A)}[thm]{\rAEb{13}}}
 \proofstepwidestar[15]{x\in B_j}{\rA}
 \proofstepwidestar[11]{B_j\subseteq U}{\FormulaRefAuto{ZeroTransportUnionContainsMember}[thm]{11}}
 \proofstepwidestar[11,15]{x\in U}{\FormulaRefAuto{A\subseteq B\dsep x\in A\vdash x\in B}[thm]{16,15}}
 \proofstepwidestar[11,15]{U\neq\varnothing}{\FormulaRefAuto{a\in A\vdash A\neq\varnothing}[thm]{17}}
 \proofstepwidestar[2,11]{U\neq\varnothing}{\rEE{14,15,18}}
 \proofstepwidestar[1,2]{U\neq\varnothing}{\rEE{10,11,19}}
 \proofstepwidestar[1,2]{U\in\PowNE S}{\FormulaRefAuto{NonemptyPowersetMembership}[thm]{\rAI{9,20}}}
\closeproofpartwide
\end{tabproofsplitwide}

\Needspace{12\baselineskip}
\FormulaThmDeltaK[Mengenprodukte verteilen sich über Vereinigungen]{%
 \begin{aligned}[t]
 G_S\dsep A\in K_S\dsep J\neq\varnothing\dsep
 \forall j\in J\;(B_j\in K_S)\vdash{}\\[-2pt]
 A\Bigl(\bigcup_{j\in J}B_j\Bigr)=\bigcup_{j\in J}AB_j
 \end{aligned}
}{ZeroTransportProductUnion}{%
 \DeltaRow{Mengenfamilie}{(B_j)_{j\in J}}
 \DeltaRow{Mengen}{S\dsep J\dsep A\dsep j\dsep a\dsep b\dsep x}}
Hier seien \(U=\bigcup_{j\in J}B_j\) und
\(V=\bigcup_{j\in J}AB_j\).
\begin{tabproofsplitwide}
\proofpartwide{Jedes Produkt mit der Vereinigung liegt in einem Teilprodukt}
 \proofstepwidestar[1]{G_S}{\rA}
 \proofstepwidestar[2]{A\in K_S}{\rA}
 \proofstepwidestar[3]{J\neq\varnothing}{\rA}
 \proofstepwidestar[4]{\forall j\in J\;(B_j\in K_S)}{\rA}
 \proofstepwidestar[3,4]{U\in K_S}{\FormulaRefAuto{ZeroTransportUnionCarrier}[thm]{3,4}}
 \proofstepwidestar[6]{x\in AU}{\rA}
 \proofstepwidestar[1,2,3,4,6]{\exists a\in A\,\exists b\in U\;(x=ab)}{\begin{gathered}
      \RuleRef{AE2}{\land E2}(\RuleRef{RE}{\rightarrow E}(\\[-2pt]
      \RuleRef{LRE1}{\leftrightarrow E1}(\\[-2pt]
      \FormulaRefAuto{PowerSemigroupProductMembership}[thm](\\[-2pt]
      1,2,5)),6))
    \end{gathered}}
 \proofstepwidestar[8]{a\in A\land(b\in U\land x=ab)}{\rA}
 \proofstepwidestar[8]{\exists j\in J\;(b\in B_j)}{\FormulaRefAuto{SemigroupIndexedUnionNotationDef}[def]{\rAEa{\rAEb{8}}}}
 \proofstepwidestar[10]{j\in J\land b\in B_j}{\rA}
 \proofstepwidestar[4,10]{B_j\in K_S}{\rRE{\rUE{4},\rAEa{10}}}
 \proofstepwidestar[1,2,4,8,10]{ab\in AB_j}{\FormulaRefAuto{PowerSemigroupProductContains}[thm]{1,2,11,\rAEa{8},\rAEb{10}}}
 \proofstepwidestar[1,2,4,8,10]{x\in AB_j}{\rIE{\FormulaRefAuto{a=b\vdash b=a}[thm]{\rAEb{\rAEb{8}}},12}}
 \proofstepwidestar[1,2,4,8,10]{x\in V}{\FormulaRefAuto{SemigroupIndexedUnionNotationDef}[def]{\rEI{\rAI{\rAEa{10},13}}}}
 \proofstepwidestar[1,2,4,8]{x\in V}{\rEE{9,10,14}}
 \proofstepwidestar[1,2,3,4,6]{x\in V}{\rEEStar{7,8,15}}
 \proofstepwidestar[1,2,3,4]{AU\subseteq V}{\FormulaRefAuto{A\subseteq B\coloneq\forall x\,(x\in A\rightarrow x\in B)}[def]{\rUI{\rRI{6,16}}}}
\closeproofpartwide
\proofpartwide{Jedes Teilprodukt liegt im Produkt mit der Vereinigung}
 \setcounter{proofstepnr}{17}
 \proofstepwidestar[18]{x\in V}{\rA}
 \proofstepwidestar[18]{\exists j\in J\;(x\in AB_j)}{\FormulaRefAuto{SemigroupIndexedUnionNotationDef}[def]{18}}
 \proofstepwidestar[20]{j\in J\land x\in AB_j}{\rA}
 \proofstepwidestar[4,20]{B_j\in K_S}{\rRE{\rUE{4},\rAEa{20}}}
 \proofstepwidestar[1,2,4,20]{\exists a\in A\,\exists b\in B_j\;(x=ab)}{\begin{gathered}
      \RuleRef{AE2}{\land E2}(\RuleRef{RE}{\rightarrow E}(\\[-2pt]
      \RuleRef{LRE1}{\leftrightarrow E1}(\\[-2pt]
      \FormulaRefAuto{PowerSemigroupProductMembership}[thm](\\[-2pt]
      1,2,21)),\rAEb{20}))
    \end{gathered}}
 \proofstepwidestar[23]{a\in A\land(b\in B_j\land x=ab)}{\rA}
 \proofstepwidestar[20]{B_j\subseteq U}{\FormulaRefAuto{ZeroTransportUnionContainsMember}[thm]{\rAEa{20}}}
 \proofstepwidestar[20,23]{b\in U}{\FormulaRefAuto{A\subseteq B\dsep x\in A\vdash x\in B}[thm]{24,\rAEa{\rAEb{23}}}}
 \proofstepwidestar[1,2,3,4,20,23]{ab\in AU}{\FormulaRefAuto{PowerSemigroupProductContains}[thm]{1,2,5,\rAEa{23},25}}
 \proofstepwidestar[1,2,3,4,20,23]{x\in AU}{\rIE{\FormulaRefAuto{a=b\vdash b=a}[thm]{\rAEb{\rAEb{23}}},26}}
 \proofstepwidestar[1,2,3,4,20]{x\in AU}{\rEEStar{22,23,27}}
 \proofstepwidestar[1,2,3,4,18]{x\in AU}{\rEE{19,20,28}}
 \proofstepwidestar[1,2,3,4]{V\subseteq AU}{\FormulaRefAuto{A\subseteq B\coloneq\forall x\,(x\in A\rightarrow x\in B)}[def]{\rUI{\rRI{18,29}}}}
 \proofstepwidestar[1,2,3,4]{AU=V}{\FormulaRefAuto{A\subseteq B\dsep B\subseteq A\vdash A=B}[thm]{17,30}}
\closeproofpartwide
\end{tabproofsplitwide}

\Needspace{12\baselineskip}
\FormulaThmDeltaK[Eine surjektive Umindizierung ändert die Vereinigung nicht]{%
 \begin{aligned}[t]
 &\forall i\in J\;(r(i)\in L)\dsep
 \forall j\in L\,\exists i\in J\;(r(i)=j)\vdash{}\\[-2pt]
 &\bigcup_{i\in J}B_{r(i)}=\bigcup_{j\in L}B_j
 \end{aligned}
}{ZeroTransportReindexedUnion}{%
 \DeltaRow{Mengenfamilie}{(B_j)_{j\in L}}
 \DeltaRow{Mengenterm}{r(i)}
 \DeltaRow{Mengen}{J\dsep L\dsep i\dsep j\dsep x}}
Hier seien \(U=\bigcup_{i\in J}B_{r(i)}\) und
\(V=\bigcup_{j\in L}B_j\).
\begin{tabproofsplitwide}
\proofpartwide{Keine neuen Indizes}
 \proofstepwidestar[1]{\forall i\in J\;(r(i)\in L)}{\rA}
 \proofstepwidestar[2]{\forall j\in L\,\exists i\in J\;(r(i)=j)}{\rA}
 \proofstepwidestar[3]{x\in U}{\rA}
 \proofstepwidestar[3]{\exists i\in J\;(x\in B_{r(i)})}{\FormulaRefAuto{SemigroupIndexedUnionNotationDef}[def]{3}}
 \proofstepwidestar[5]{i\in J\land x\in B_{r(i)}}{\rA}
 \proofstepwidestar[1,5]{r(i)\in L}{\rRE{\rUE{1},\rAEa{5}}}
 \proofstepwidestar[1,5]{x\in V}{\FormulaRefAuto{SemigroupIndexedUnionNotationDef}[def]{\rEI{\rAI{6,\rAEb{5}}}}}
 \proofstepwidestar[1,3]{x\in V}{\rEE{4,5,7}}
 \proofstepwidestar[1]{U\subseteq V}{\FormulaRefAuto{A\subseteq B\coloneq\forall x\,(x\in A\rightarrow x\in B)}[def]{\rUI{\rRI{3,8}}}}
\closeproofpartwide
\proofpartwide{Jeder ursprüngliche Index wird getroffen}
 \setcounter{proofstepnr}{9}
 \proofstepwidestar[10]{x\in V}{\rA}
 \proofstepwidestar[10]{\exists j\in L\;(x\in B_j)}{\FormulaRefAuto{SemigroupIndexedUnionNotationDef}[def]{10}}
 \proofstepwidestar[12]{j\in L\land x\in B_j}{\rA}
 \proofstepwidestar[2,12]{\exists i\in J\;(r(i)=j)}{\rRE{\rUE{2},\rAEa{12}}}
 \proofstepwidestar[14]{i\in J\land r(i)=j}{\rA}
 \proofstepwidestar[12,14]{x\in B_{r(i)}}{\rIE{\FormulaRefAuto{a=b\vdash b=a}[thm]{\rAEb{14}},\rAEb{12}}}
 \proofstepwidestar[12,14]{x\in U}{\FormulaRefAuto{SemigroupIndexedUnionNotationDef}[def]{\rEI{\rAI{\rAEa{14},15}}}}
 \proofstepwidestar[2,12]{x\in U}{\rEE{13,14,16}}
 \proofstepwidestar[2,10]{x\in U}{\rEE{11,12,17}}
 \proofstepwidestar[2]{V\subseteq U}{\FormulaRefAuto{A\subseteq B\coloneq\forall x\,(x\in A\rightarrow x\in B)}[def]{\rUI{\rRI{10,18}}}}
 \proofstepwidestar[1,2]{U=V}{\FormulaRefAuto{A\subseteq B\dsep B\subseteq A\vdash A=B}[thm]{9,19}}
\closeproofpartwide
\end{tabproofsplitwide}

\Needspace{12\baselineskip}
\FormulaThmDeltaK[Punktweise gleiche Mengen haben dieselbe Vereinigung]{%
 \forall i\in J\;(B_i=C_i)\vdash
 \bigcup_{i\in J}B_i=\bigcup_{i\in J}C_i
}{ZeroTransportUnionPointwise}{%
 \DeltaRow{Mengenfamilien}{(B_i)_{i\in J}\dsep(C_i)_{i\in J}}
 \DeltaRow{Mengen}{J\dsep i\dsep x}}
Hier seien \(U=\bigcup_{i\in J}B_i\) und
\(V=\bigcup_{i\in J}C_i\).
\begin{tabproofsplitwide}
\proofpartwide{Die erste Vereinigung liegt in der zweiten}
 \proofstepwidestar[1]{\forall i\in J\;(B_i=C_i)}{\rA}
 \proofstepwidestar[2]{x\in U}{\rA}
 \proofstepwidestar[2]{\exists i\in J\;(x\in B_i)}{\FormulaRefAuto{SemigroupIndexedUnionNotationDef}[def]{2}}
 \proofstepwidestar[4]{i\in J\land x\in B_i}{\rA}
 \proofstepwidestar[1,4]{B_i=C_i}{\rRE{\rUE{1},\rAEa{4}}}
 \proofstepwidestar[1,4]{x\in C_i}{\rIE{5,\rAEb{4}}}
 \proofstepwidestar[1,4]{x\in V}{\FormulaRefAuto{SemigroupIndexedUnionNotationDef}[def]{\rEI{\rAI{\rAEa{4},6}}}}
 \proofstepwidestar[1,2]{x\in V}{\rEE{3,4,7}}
 \proofstepwidestar[1]{U\subseteq V}{\FormulaRefAuto{A\subseteq B\coloneq\forall x\,(x\in A\rightarrow x\in B)}[def]{\rUI{\rRI{2,8}}}}
\closeproofpartwide
\proofpartwide{Die zweite Vereinigung liegt in der ersten}
 \setcounter{proofstepnr}{9}
 \proofstepwidestar[10]{x\in V}{\rA}
 \proofstepwidestar[10]{\exists i\in J\;(x\in C_i)}{\FormulaRefAuto{SemigroupIndexedUnionNotationDef}[def]{10}}
 \proofstepwidestar[12]{i\in J\land x\in C_i}{\rA}
 \proofstepwidestar[1,12]{B_i=C_i}{\rRE{\rUE{1},\rAEa{12}}}
 \proofstepwidestar[1,12]{C_i=B_i}{\FormulaRefAuto{a=b\vdash b=a}[thm]{13}}
 \proofstepwidestar[1,12]{x\in B_i}{\rIE{14,\rAEb{12}}}
 \proofstepwidestar[1,12]{x\in U}{\FormulaRefAuto{SemigroupIndexedUnionNotationDef}[def]{\rEI{\rAI{\rAEa{12},15}}}}
 \proofstepwidestar[1,10]{x\in U}{\rEE{11,12,16}}
 \proofstepwidestar[1]{V\subseteq U}{\FormulaRefAuto{A\subseteq B\coloneq\forall x\,(x\in A\rightarrow x\in B)}[def]{\rUI{\rRI{10,17}}}}
 \proofstepwidestar[1]{U=V}{\FormulaRefAuto{A\subseteq B\dsep B\subseteq A\vdash A=B}[thm]{9,18}}
\closeproofpartwide
\end{tabproofsplitwide}

\Needspace{12\baselineskip}
\FormulaThmDeltaK[Gemeinsame Wirkung der Einermengen]{%
 G_T\dsep D,C\in K_T\dsep\forall t\in D\;(\{t\}C=C)
 \vdash DC=C
}{ZeroTransportSingletonsFixUnion}{%
 \DeltaRow{Mengen}{T\dsep D\dsep C\dsep t\dsep c\dsep x\dsep u}}
\begin{tabproofsplitwide}
\proofpartwide{Das gesamte Produkt bleibt in der festen Menge}
 \proofstepwidestar[1]{G_T}{\rA}
 \proofstepwidestar[2]{D,C\in K_T}{\rA}
 \proofstepwidestar[3]{\forall t\in D\;(\{t\}C=C)}{\rA}
 \proofstepwidestar[4]{x\in DC}{\rA}
 \proofstepwidestar[1,2,4]{\exists t\in D\,\exists c\in C\;(x=tc)}{\begin{gathered}
      \RuleRef{AE2}{\land E2}(\RuleRef{RE}{\rightarrow E}(\\[-2pt]
      \RuleRef{LRE1}{\leftrightarrow E1}(\\[-2pt]
      \FormulaRefAuto{PowerSemigroupProductMembership}[thm](\\[-2pt]
      1,\rAEa{2},\\[-2pt]
      \rAEb{2})),4))
    \end{gathered}}
 \proofstepwidestar[6]{t\in D\land(c\in C\land x=tc)}{\rA}
 \proofstepwidestar[2,6]{t\in T}{\FormulaRefAuto{PowerSemigroupCarrierElement}[thm]{2,\rAEa{6}}}
 \proofstepwidestar[2,6]{\{t\}\in K_T}{\FormulaRefAuto{SingletonPowerMember}[thm]{7}}
 \proofstepwidestar[]{t\in\{t\}}{\FormulaRefAuto{a\in\{a\}}[thm]}
 \proofstepwidestar[1,2,6]{tc\in\{t\}C}{\FormulaRefAuto{PowerSemigroupProductContains}[thm]{1,8,2,9,\rAEa{\rAEb{6}}}}
 \proofstepwidestar[3,6]{\{t\}C=C}{\rRE{\rUE{3},\rAEa{6}}}
 \proofstepwidestar[1,2,3,6]{x\in C}{\rIE{11,\rIE{\FormulaRefAuto{a=b\vdash b=a}[thm]{\rAEb{\rAEb{6}}},10}}}
 \proofstepwidestar[1,2,3,4]{x\in C}{\rEEStar{5,6,12}}
 \proofstepwidestar[1,2,3]{DC\subseteq C}{\FormulaRefAuto{A\subseteq B\coloneq\forall x\,(x\in A\rightarrow x\in B)}[def]{\rUI{\rRI{4,13}}}}
\closeproofpartwide
\proofpartwide{Schon eine Einermenge trifft die ganze feste Menge}
 \setcounter{proofstepnr}{14}
 \proofstepwidestar[2]{D\subseteq T\land D\neq\varnothing}{\FormulaRefAuto{NonemptyPowersetMembership}[thm]{2}}
 \proofstepwidestar[2]{\exists t\;(t\in D)}{\FormulaRefAuto{A\neq\varnothing\eqvdash\exists x\,(x\in A)}[thm]{\rAEb{15}}}
 \proofstepwidestar[17]{t\in D}{\rA}
 \proofstepwidestar[2,17]{t\in T}{\FormulaRefAuto{PowerSemigroupCarrierElement}[thm]{2,17}}
 \proofstepwidestar[2,17]{\{t\}\in K_T}{\FormulaRefAuto{SingletonPowerMember}[thm]{18}}
 \proofstepwidestar[3,17]{\{t\}C=C}{\rRE{\rUE{3},17}}
 \proofstepwidestar[21]{x\in C}{\rA}
 \proofstepwidestar[3,17,21]{x\in\{t\}C}{\rIE{\FormulaRefAuto{a=b\vdash b=a}[thm]{20},21}}
 \proofstepwidestar[1,2,3,17,21]{\exists u\in\{t\}\,\exists c\in C\;(x=uc)}{\begin{gathered}
      \RuleRef{AE2}{\land E2}(\RuleRef{RE}{\rightarrow E}(\\[-2pt]
      \RuleRef{LRE1}{\leftrightarrow E1}(\\[-2pt]
      \FormulaRefAuto{PowerSemigroupProductMembership}[thm](\\[-2pt]
      1,19,\rAEb{2})),22))
    \end{gathered}}
 \proofstepwidestar[24]{u\in\{t\}\land(c\in C\land x=uc)}{\rA}
 \proofstepwidestar[24]{u=t}{\FormulaRefAuto{x\in\{a\}\eqvdash x=a}[thm]{\rAEa{24}}}
 \proofstepwidestar[17,24]{u\in D}{\rIE{\FormulaRefAuto{a=b\vdash b=a}[thm]{25},17}}
 \proofstepwidestar[1,2,17,24]{uc\in DC}{\FormulaRefAuto{PowerSemigroupProductContains}[thm]{1,2,26,\rAEa{\rAEb{24}}}}
 \proofstepwidestar[1,2,17,24]{x\in DC}{\rIE{\FormulaRefAuto{a=b\vdash b=a}[thm]{\rAEb{\rAEb{24}}},27}}
 \proofstepwidestar[1,2,3,17,21]{x\in DC}{\rEEStar{23,24,28}}
 \proofstepwidestar[1,2,3,21]{x\in DC}{\rEE{16,17,29}}
 \proofstepwidestar[1,2,3]{C\subseteq DC}{\FormulaRefAuto{A\subseteq B\coloneq\forall x\,(x\in A\rightarrow x\in B)}[def]{\rUI{\rRI{21,30}}}}
 \proofstepwidestar[1,2,3]{DC=C}{\FormulaRefAuto{A\subseteq B\dsep B\subseteq A\vdash A=B}[thm]{14,31}}
\closeproofpartwide
\end{tabproofsplitwide}

\section{Potenznull und Urbilder der Einermengen}

\Needspace{12\baselineskip}
\FormulaThmDeltaK[Das Bild der Null ist eine Null der Potenzhalbgruppe]{%
 H_S\dsep I\vdash\SemigroupWithZeroAlg{K_T}{\diamond_{\mathcal P}}{Z}
}{ZeroTransportPowerZero}{\DeltaRow{Mengen}{S\dsep T\dsep0_S\dsep X\dsep Y}}
\begin{tabproofwide}
 \proofstepwidestar[1]{H_S}{\rA}
 \proofstepwidestar[2]{I}{\rA}
 \proofstepwidestar[1]{\SemigroupWithZeroAlg{K_S}{\star_{\mathcal P}}{z_S}}{\FormulaRefAuto{PowerSemigroupWithZero}[thm]{1}}
 \proofstepwidestar[1]{z_S\in K_S}{\FormulaRefAuto{SemigroupWithZeroCarrierAxiom}[ax]{3}}
 \proofstepwidestar[2]{\Phi:K_S\to K_T}{\FormulaRefAuto{SemigroupIsoFunction}[thm]{2}}
 \proofstepwidestar[1,2]{Z\in K_T}{\FormulaRefAuto{F\colon A\to B\dsep x\in A\vdash F(x)\in B}[thm]{5,4}}
 \proofstepwidestar[1]{\forall X\in K_S\;(z_SX=z_S\land Xz_S=z_S)}{\FormulaRefAuto{SemigroupWithZeroDef}[def]{3}}
 \proofstepwidestar[1,2]{\begin{aligned}[t]
 &(\forall X\in K_S\;(z_SX=z_S\land Xz_S=z_S))\\[-2pt]
 &\quad\leftrightarrow(\forall Y\in K_T\;(ZY=Z\land YZ=Z))
 \end{aligned}}{\FormulaRefAuto{SemigroupIsoIdentityZeroConditions}[thm]{2,\rAI{4,4}}}
 \proofstepwidestar[1,2]{\forall Y\in K_T\;(ZY=Z\land YZ=Z)}{\rRE{\rLREa{8},7}}
 \proofstepwidestar[2]{\SemigroupAlg{K_T}{\diamond_{\mathcal P}}}{\FormulaRefAuto{SemigroupIsoTarget}[thm]{2}}
 \proofstepwidestar[1,2]{\SemigroupWithZeroAlg{K_T}{\diamond_{\mathcal P}}{Z}}{\FormulaRefAuto{SemigroupWithZeroDef}[def]{10,6,9}}
\end{tabproofwide}

\Needspace{12\baselineskip}
\FormulaThmDeltaK[Die Urbilder der Einermengen sind nichtleere Teilmengen]{%
 I\dsep t\in T\vdash A_t\in K_S\land\Phi(A_t)=\{t\}
}{ZeroTransportInverseSingletons}{\DeltaRow{Mengen}{S\dsep T\dsep t}}
\begin{tabproofwide}
 \proofstepwidestar[1]{I}{\rA}
 \proofstepwidestar[2]{t\in T}{\rA}
 \proofstepwidestar[2]{\{t\}\in K_T}{\FormulaRefAuto{SingletonPowerMember}[thm]{2}}
 \proofstepwidestar[1]{\Phi^{-1}:K_T\bij K_S}{\FormulaRefAuto{SemigroupIsoInverseCarrier}[thm]{1}}
 \proofstepwidestar[1]{\Phi^{-1}:K_T\to K_S}{\FormulaRefAuto{Bijektive Funktion}[def]{4}}
 \proofstepwidestar[1,2]{A_t\in K_S}{\FormulaRefAuto{F\colon A\to B\dsep x\in A\vdash F(x)\in B}[thm]{5,3}}
 \proofstepwidestar[1]{\Phi:K_S\bij K_T}{\FormulaRefAuto{SemigroupIsoBijectiveAxiom}[ax]{1}}
 \proofstepwidestar[1,2]{\Phi(A_t)=\{t\}}{\FormulaRefAuto{InverseFunctionRightInverse}[thm]{7,3}}
 \proofstepwidestar[1,2]{A_t\in K_S\land\Phi(A_t)=\{t\}}{\rAI{6,8}}
\end{tabproofwide}

\Needspace{12\baselineskip}
\FormulaThmDeltaK[Produkte der Urbilder von Einermengen]{%
 G_T\dsep I\dsep u,v\in T\vdash A_uA_v=A_{uv}
}{ZeroTransportInverseSingletonProduct}{\DeltaRow{Mengen}{S\dsep T\dsep u\dsep v}}
\begin{tabproofwide}
 \proofstepwidestar[1]{G_T}{\rA}
 \proofstepwidestar[2]{I}{\rA}
 \proofstepwidestar[3]{u,v\in T}{\rA}
 \proofstepwidestar[3]{\{u\},\{v\}\in K_T}{\FormulaRefAuto{SingletonPowerMember}[thm]{3}}
 \proofstepwidestar[1,3]{\{u\}\{v\}=\{uv\}}{\FormulaRefAuto{SingletonPowerProduct}[thm]{1,3}}
 \proofstepwidestar[2,3]{\Phi^{-1}(\{u\}\{v\})=A_uA_v}{\FormulaRefAuto{SemigroupIsoInverseOperation}[thm]{2,4}}
 \proofstepwidestar[1,2,3]{A_{uv}=A_uA_v}{\rIE{5,6}}
 \proofstepwidestar[1,2,3]{A_uA_v=A_{uv}}{\FormulaRefAuto{a=b\vdash b=a}[thm]{7}}
\end{tabproofwide}

\Needspace{12\baselineskip}
\FormulaThmDeltaK[Die Vereinigung aller relevanten Urbilder ist zulässig]{%
 H_S\dsep I\vdash
 Z\in K_T\land\bigl(W\in K_S\land\forall z\in Z\;(A_z\in K_S)\bigr)
}{ZeroTransportPreimageUnionCarrier}{\DeltaRow{Mengen}{S\dsep T\dsep0_S\dsep z}}
\begin{tabproofwide}
 \proofstepwidestar[1]{H_S}{\rA}
 \proofstepwidestar[2]{I}{\rA}
 \proofstepwidestar[1,2]{\SemigroupWithZeroAlg{K_T}{\diamond_{\mathcal P}}{Z}}{\FormulaRefAuto{ZeroTransportPowerZero}[thm]{1,2}}
 \proofstepwidestar[1,2]{Z\in K_T}{\FormulaRefAuto{SemigroupWithZeroCarrierAxiom}[ax]{3}}
 \proofstepwidestar[1,2]{Z\subseteq T\land Z\neq\varnothing}{\FormulaRefAuto{NonemptyPowersetMembership}[thm]{4}}
 \proofstepwidestar[6]{z\in Z}{\rA}
 \proofstepwidestar[1,2,6]{z\in T}{\FormulaRefAuto{PowerSemigroupCarrierElement}[thm]{4,6}}
 \proofstepwidestar[1,2,6]{A_z\in K_S\land\Phi(A_z)=\{z\}}{\FormulaRefAuto{ZeroTransportInverseSingletons}[thm]{2,7}}
 \proofstepwidestar[1,2]{\forall z\in Z\;(A_z\in K_S)}{\rUI{\rRI{6,\rAEa{8}}}}
 \proofstepwidestar[1,2]{W\in K_S}{\FormulaRefAuto{ZeroTransportUnionCarrier}[thm]{\rAEb{5},9}}
 \proofstepwidestar[1,2]{Z\in K_T\land\bigl(W\in K_S\land\forall z\in Z\;(A_z\in K_S)\bigr)}{\rAI{4,\rAI{10,9}}}
\end{tabproofwide}

\section{Die entscheidende Vereinigungsrechnung}

In den nächsten beiden Hilfssätzen steht
\(\Gamma=H_S\land(G_T\land I)\) für die drei Voraussetzungen.

\Needspace{12\baselineskip}
\FormulaThmDeltaK[Jedes Urbild lässt die Vereinigung fest]{%
 \Gamma\vdash\forall z\in Z\;(A_zW=W)
}{ZeroTransportStablePreimageUnion}{%
 \DeltaRow{Mengen}{S\dsep T\dsep0_S\dsep z\dsep w\dsep v\dsep u}}
\begin{tabproofsplitwide}
\proofpartwide{Multiplikation mit einem Index bleibt in der Indexmenge}
 \proofstepwidestar[1]{\Gamma}{\rA}
 \proofstepwidestar[1]{H_S}{\rAEa{1}}
 \proofstepwidestar[1]{G_T}{\rAEa{\rAEb{1}}}
 \proofstepwidestar[1]{I}{\rAEb{\rAEb{1}}}
 \proofstepwidestar[1]{G_S}{\FormulaRefAuto{SemigroupWithZeroBaseAxiom}[ax]{2}}
 \proofstepwidestar[1]{\SemigroupWithZeroAlg{K_T}{\diamond_{\mathcal P}}{Z}}{\FormulaRefAuto{ZeroTransportPowerZero}[thm]{2,4}}
 \proofstepwidestar[1]{Z\in K_T\land\bigl(W\in K_S\land\forall w\in Z\;(A_w\in K_S)\bigr)}{\FormulaRefAuto{ZeroTransportPreimageUnionCarrier}[thm]{2,4}}
 \proofstepwidestar[1]{Z\subseteq T\land Z\neq\varnothing}{\FormulaRefAuto{NonemptyPowersetMembership}[thm]{\rAEa{7}}}
 \proofstepwidestar[9]{z\in Z}{\rA}
 \proofstepwidestar[1,9]{z\in T}{\FormulaRefAuto{PowerSemigroupCarrierElement}[thm]{\rAEa{7},9}}
 \proofstepwidestar[1,9]{A_z\in K_S\land\Phi(A_z)=\{z\}}{\FormulaRefAuto{ZeroTransportInverseSingletons}[thm]{4,10}}
 \proofstepwidestar[1,9]{\{z\}\in K_T}{\FormulaRefAuto{SingletonPowerMember}[thm]{10}}
 \proofstepwidestar[1,9]{\{z\}Z=Z}{\FormulaRefAuto{SemigroupWithZeroRightAbsorptionAxiom}[ax]{6,12}}
 \proofstepwidestar[14]{w\in Z}{\rA}
 \proofstepwidestar[1,14]{w\in T}{\FormulaRefAuto{PowerSemigroupCarrierElement}[thm]{\rAEa{7},14}}
 \proofstepwidestar[]{z\in\{z\}}{\FormulaRefAuto{a\in\{a\}}[thm]}
 \proofstepwidestar[1,9,14]{zw\in\{z\}Z}{\FormulaRefAuto{PowerSemigroupProductContains}[thm]{3,12,\rAEa{7},16,14}}
 \proofstepwidestar[1,9,14]{zw\in Z}{\rIE{13,17}}
 \proofstepwidestar[1,9]{\forall w\in Z\;(zw\in Z)}{\rUI{\rRI{14,18}}}
\closeproofpartwide
\proofpartwide{Jeder Index wird getroffen}
 \setcounter{proofstepnr}{19}
 \proofstepwidestar[20]{v\in Z}{\rA}
 \proofstepwidestar[1,9,20]{v\in\{z\}Z}{\rIE{\FormulaRefAuto{a=b\vdash b=a}[thm]{13},20}}
 \proofstepwidestar[1,9,20]{\exists u\in\{z\}\,\exists w\in Z\;(v=uw)}{\begin{gathered}
      \RuleRef{AE2}{\land E2}(\RuleRef{RE}{\rightarrow E}(\\[-2pt]
      \RuleRef{LRE1}{\leftrightarrow E1}(\\[-2pt]
      \FormulaRefAuto{PowerSemigroupProductMembership}[thm](\\[-2pt]
      3,12,\rAEa{7})),21))
    \end{gathered}}
 \proofstepwidestar[23]{u\in\{z\}\land(w\in Z\land v=uw)}{\rA}
 \proofstepwidestar[23]{u=z}{\FormulaRefAuto{x\in\{a\}\eqvdash x=a}[thm]{\rAEa{23}}}
 \proofstepwidestar[23]{v=zw}{\rIE{24,\rAEb{\rAEb{23}}}}
 \proofstepwidestar[23]{zw=v}{\FormulaRefAuto{a=b\vdash b=a}[thm]{25}}
 \proofstepwidestar[23]{\exists w\in Z\;(zw=v)}{\rEI{\rAI{\rAEa{\rAEb{23}},26}}}
 \proofstepwidestar[1,9,20]{\exists w\in Z\;(zw=v)}{\rEEStar{22,23,27}}
 \proofstepwidestar[1,9]{\forall v\in Z\,\exists w\in Z\;(zw=v)}{\rUI{\rRI{20,28}}}
 \proofstepwidestar[1,9]{\bigcup_{w\in Z}A_{zw}=W}{\FormulaRefAuto{ZeroTransportReindexedUnion}[thm]{19,29}}
\closeproofpartwide
\Needspace{8\baselineskip}
\proofpartwide{Distributivität und Umindizierung}
 \setcounter{proofstepnr}{30}
 \proofstepwidestar[1,9,14]{A_zA_w=A_{zw}}{\FormulaRefAuto{ZeroTransportInverseSingletonProduct}[thm]{3,4,\rAI{10,15}}}
 \proofstepwidestar[1,9]{\forall w\in Z\;(A_zA_w=A_{zw})}{\rUI{\rRI{14,31}}}
 \proofstepwidestar[1,9]{\bigcup_{w\in Z}A_zA_w=\bigcup_{w\in Z}A_{zw}}{\FormulaRefAuto{ZeroTransportUnionPointwise}[thm]{32}}
 \proofstepwidestar[1,9]{A_zW=\bigcup_{w\in Z}A_zA_w}{\FormulaRefAuto{ZeroTransportProductUnion}[thm]{5,\rAEa{11},\rAEb{8},\rAEb{\rAEb{7}}}}
 \proofstepwidestar[1,9]{A_zW=W}{\rChain{34,33,30}}
 \proofstepwidestar[1]{\forall z\in Z\;(A_zW=W)}{\rUI{\rRI{9,35}}}
\closeproofpartwide
\end{tabproofsplitwide}

\Needspace{12\baselineskip}
\FormulaThmDeltaK[Die Vereinigung ist die ursprüngliche Null-Einermenge]{%
 \Gamma\vdash W=z_S
}{ZeroTransportPreimageUnionIsZero}{\DeltaRow{Mengen}{S\dsep T\dsep0_S\dsep z}}
\begin{tabproofsplitwide}
\proofpartwide{Die Einermengen der Potenznull lassen das Bild fest}
 \proofstepwidestar[1]{\Gamma}{\rA}
 \proofstepwidestar[1]{H_S}{\rAEa{1}}
 \proofstepwidestar[1]{G_T}{\rAEa{\rAEb{1}}}
 \proofstepwidestar[1]{I}{\rAEb{\rAEb{1}}}
 \proofstepwidestar[1]{\SemigroupWithZeroAlg{K_T}{\diamond_{\mathcal P}}{Z}}{\FormulaRefAuto{ZeroTransportPowerZero}[thm]{2,4}}
 \proofstepwidestar[1]{Z\in K_T\land\bigl(W\in K_S\land\forall z\in Z\;(A_z\in K_S)\bigr)}{\FormulaRefAuto{ZeroTransportPreimageUnionCarrier}[thm]{2,4}}
 \proofstepwidestar[1]{\Phi:K_S\to K_T}{\FormulaRefAuto{SemigroupIsoFunction}[thm]{4}}
 \proofstepwidestar[1]{\Phi(W)\in K_T}{\FormulaRefAuto{F\colon A\to B\dsep x\in A\vdash F(x)\in B}[thm]{7,\rAEa{\rAEb{6}}}}
 \proofstepwidestar[1]{\forall z\in Z\;(A_zW=W)}{\FormulaRefAuto{ZeroTransportStablePreimageUnion}[thm]{1}}
 \proofstepwidestar[10]{z\in Z}{\rA}
 \proofstepwidestar[1,10]{z\in T}{\FormulaRefAuto{PowerSemigroupCarrierElement}[thm]{\rAEa{6},10}}
 \proofstepwidestar[1,10]{A_z\in K_S\land\Phi(A_z)=\{z\}}{\FormulaRefAuto{ZeroTransportInverseSingletons}[thm]{4,11}}
 \proofstepwidestar[1,10]{A_zW=W}{\rRE{\rUE{9},10}}
 \proofstepwidestar[1,10]{\Phi(A_zW)=\Phi(A_z)\Phi(W)}{\FormulaRefAuto{SemigroupIsoOperation}[thm]{4,\rAI{\rAEa{12},\rAEa{\rAEb{6}}}}}
 \proofstepwidestar[1,10]{\Phi(W)=\Phi(A_z)\Phi(W)}{\rIE{13,14}}
 \proofstepwidestar[1,10]{\Phi(W)=\{z\}\Phi(W)}{\rIE{\rAEb{12},15}}
 \proofstepwidestar[1,10]{\{z\}\Phi(W)=\Phi(W)}{\FormulaRefAuto{a=b\vdash b=a}[thm]{16}}
 \proofstepwidestar[1]{\forall z\in Z\;(\{z\}\Phi(W)=\Phi(W))}{\rUI{\rRI{10,17}}}
\closeproofpartwide
\proofpartwide{Nullwirkung und Injektivität}
 \setcounter{proofstepnr}{18}
 \proofstepwidestar[1]{Z\Phi(W)=\Phi(W)}{\FormulaRefAuto{ZeroTransportSingletonsFixUnion}[thm]{3,\rAI{\rAEa{6},8},18}}
 \proofstepwidestar[1]{Z\Phi(W)=Z}{\FormulaRefAuto{SemigroupWithZeroLeftAbsorptionAxiom}[ax]{5,8}}
 \proofstepwidestar[1]{\Phi(W)=Z}{\rChain{\FormulaRefAuto{a=b\vdash b=a}[thm]{19},20}}
 \proofstepwidestar[1]{\SemigroupWithZeroAlg{K_S}{\star_{\mathcal P}}{z_S}}{\FormulaRefAuto{PowerSemigroupWithZero}[thm]{2}}
 \proofstepwidestar[1]{z_S\in K_S}{\FormulaRefAuto{SemigroupWithZeroCarrierAxiom}[ax]{22}}
 \proofstepwidestar[1]{W=z_S\leftrightarrow\Phi(W)=\Phi(z_S)}{\FormulaRefAuto{SemigroupIsoEqualityReflection}[thm]{4,\rAI{\rAEa{\rAEb{6}},23}}}
 \proofstepwidestar[1]{W=z_S}{\rRE{\rLREb{24},21}}
\closeproofpartwide
\end{tabproofsplitwide}

\section{Der allgemeine Nulltransportsatz}

\Needspace{14\baselineskip}
\FormulaThmDeltaK[Potenzisomorphismen erhalten Nullelemente]{%
 \begin{aligned}[t]
 H_S\dsep G_T\dsep I\vdash{}\\[-2pt]
 \exists 0_T\;\bigl(\SemigroupWithZeroAlg{T}{\diamond}{0_T}
 \land\Phi(\{0_S\})=\{0_T\}\bigr)
 \end{aligned}
}{PowerSemigroupZeroTransport}{%
 \DeltaRow{Mengen}{S\dsep T\dsep0_S\dsep0_T\dsep z\dsep t}
 \DeltaRow{Funktionen}{\star\dsep\diamond\dsep\Phi}}
\begin{tabproofsplitwide}
\proofpartwide{Das Bild der ursprünglichen Null ist eine Einermenge}
 \proofstepwidestar[1]{H_S}{\rA}
 \proofstepwidestar[2]{G_T}{\rA}
 \proofstepwidestar[3]{I}{\rA}
 \proofstepwidestar[1,2,3]{\Gamma}{\rAI{1,\rAI{2,3}}}
 \proofstepwidestar[1,3]{\SemigroupWithZeroAlg{K_T}{\diamond_{\mathcal P}}{Z}}{\FormulaRefAuto{ZeroTransportPowerZero}[thm]{1,3}}
 \proofstepwidestar[1,3]{Z\in K_T\land\bigl(W\in K_S\land\forall z\in Z\;(A_z\in K_S)\bigr)}{\FormulaRefAuto{ZeroTransportPreimageUnionCarrier}[thm]{1,3}}
 \proofstepwidestar[1,2,3]{W=z_S}{\FormulaRefAuto{ZeroTransportPreimageUnionIsZero}[thm]{4}}
 \proofstepwidestar[1,3]{Z\subseteq T\land Z\neq\varnothing}{\FormulaRefAuto{NonemptyPowersetMembership}[thm]{\rAEa{6}}}
 \proofstepwidestar[1,3]{\exists z\;(z\in Z)}{\FormulaRefAuto{A\neq\varnothing\eqvdash\exists x\,(x\in A)}[thm]{\rAEb{8}}}
 \proofstepwidestar[10]{z\in Z}{\rA}
 \proofstepwidestar[1,3,10]{z\in T}{\FormulaRefAuto{PowerSemigroupCarrierElement}[thm]{\rAEa{6},10}}
 \proofstepwidestar[1,3,10]{A_z\in K_S\land\Phi(A_z)=\{z\}}{\FormulaRefAuto{ZeroTransportInverseSingletons}[thm]{3,11}}
 \proofstepwidestar[10]{A_z\subseteq W}{\FormulaRefAuto{ZeroTransportUnionContainsMember}[thm]{10}}
 \proofstepwidestar[1,2,3,10]{A_z\subseteq z_S}{\rIE{7,13}}
 \proofstepwidestar[1,3,10]{A_z\subseteq S\land A_z\neq\varnothing}{\FormulaRefAuto{NonemptyPowersetMembership}[thm]{\rAEa{12}}}
 \proofstepwidestar[1,2,3,10]{A_z=z_S}{\FormulaRefAuto{NonemptySubsetSingleton}[thm]{14,\rAEb{15}}}
 \proofstepwidestar[1,2,3,10]{\Phi(A_z)=\Phi(z_S)}{\rIE{16,\rII}}
 \proofstepwidestar[1,2,3,10]{\{z\}=Z}{\rChain{\FormulaRefAuto{a=b\vdash b=a}[thm]{\rAEb{12}},17}}
 \proofstepwidestar[1,2,3,10]{Z=\{z\}}{\FormulaRefAuto{a=b\vdash b=a}[thm]{18}}
\closeproofpartwide
\proofpartwide{Ihr Grundelement absorbiert von beiden Seiten}
 \setcounter{proofstepnr}{19}
 \proofstepwidestar[20]{t\in T}{\rA}
 \proofstepwidestar[20]{\{t\}\in K_T}{\FormulaRefAuto{SingletonPowerMember}[thm]{20}}
 \proofstepwidestar[1,3,20]{Z\{t\}=Z\land\{t\}Z=Z}{\rAI{\FormulaRefAuto{SemigroupWithZeroLeftAbsorptionAxiom}[ax]{5,21},\FormulaRefAuto{SemigroupWithZeroRightAbsorptionAxiom}[ax]{5,21}}}
 \proofstepwidestar[1,2,3,10,20]{\{z\}\{t\}=\{z\}\land\{t\}\{z\}=\{z\}}{\rIE{19,22}}
 \proofstepwidestar[1,3,10]{\{z\}\in K_T}{\FormulaRefAuto{SingletonPowerMember}[thm]{11}}
 \proofstepwidestar[]{z\in\{z\}\land t\in\{t\}}{\rAI{\FormulaRefAuto{a\in\{a\}}[thm],\FormulaRefAuto{a\in\{a\}}[thm]}}
 \proofstepwidestar[1,2,3,10,20]{zt\in\{z\}\{t\}}{\FormulaRefAuto{PowerSemigroupProductContains}[thm]{2,24,21,\rAEa{25},\rAEb{25}}}
 \proofstepwidestar[1,2,3,10,20]{tz\in\{t\}\{z\}}{\FormulaRefAuto{PowerSemigroupProductContains}[thm]{2,21,24,\rAEb{25},\rAEa{25}}}
 \proofstepwidestar[1,2,3,10,20]{zt\in\{z\}}{\rIE{\rAEa{23},26}}
 \proofstepwidestar[1,2,3,10,20]{tz\in\{z\}}{\rIE{\rAEb{23},27}}
 \proofstepwidestar[1,2,3,10,20]{zt=z\land tz=z}{\rAI{\FormulaRefAuto{x\in\{a\}\eqvdash x=a}[thm]{28},\FormulaRefAuto{x\in\{a\}\eqvdash x=a}[thm]{29}}}
 \proofstepwidestar[1,2,3,10]{\forall t\in T\;(zt=z\land tz=z)}{\rUI{\rRI{20,30}}}
 \proofstepwidestar[1,2,3,10]{\SemigroupWithZeroAlg{T}{\diamond}{z}}{\FormulaRefAuto{SemigroupWithZeroDef}[def]{2,11,31}}
 \proofstepwidestar[1,2,3,10]{\exists 0_T\;\bigl(\SemigroupWithZeroAlg{T}{\diamond}{0_T}\land\Phi(\{0_S\})=\{0_T\}\bigr)}{\rEI{\rAI{32,19}}}
 \proofstepwidestar[1,2,3]{\exists 0_T\;\bigl(\SemigroupWithZeroAlg{T}{\diamond}{0_T}\land\Phi(\{0_S\})=\{0_T\}\bigr)}{\rEE{9,10,33}}
\closeproofpartwide
\end{tabproofsplitwide}

\Needspace{5\baselineskip}
Der Satz enthält zugleich die Existenz des Nullelements von \(T\) und
die genaue Gleichung für das Bild von \(\{0_S\}\).
Die Vereinigung aller Urbilder wurde innerhalb von \(S\) gebildet;
eine Erhaltung von Vereinigungen durch \(\Phi\) kommt in keiner Zeile vor.
